\documentclass[11pt,a4paper]{article}
\usepackage[margin=0.75in]{geometry}
\usepackage{graphicx}
\usepackage{hyperref}
\usepackage{tabularx}
\usepackage{booktabs}
\usepackage{float}
\usepackage{xcolor}

\title{\textbf{Task 3: Attention Visualization \& Head Clustering}}
\author{TinyLlama MLX Fine-Tuning Project}
\date{\today}

\begin{document}

\maketitle

\section*{Task Description}
This task provides \textbf{interpretability analysis} of transformer attention mechanisms through visualization and clustering. We extract attention weights from all 22 layers × 32 heads (704 total heads), analyze their patterns, and compare base vs. LoRA-tuned models to understand how fine-tuning modifies attention behavior.

\textbf{Key objectives:}
\begin{itemize}
    \item Visualize attention heatmaps for individual inference examples
    \item Cluster attention heads by similarity to identify functional groups
    \item Quantify attention pattern differences between base and fine-tuned models
\end{itemize}

\section*{Technical Implementation}

\subsection*{Procedure}
\begin{enumerate}
    \item \textbf{Attention Extraction:} Hook into transformer layers to capture attention weights during generation
    \item \textbf{Attention Rollout:} Aggregate attention across layers: $A^{(l)} = A^{(l)} \cdot A^{(l-1)}$
    \item \textbf{Head Clustering:} Compute cosine similarity, apply hierarchical clustering
    \item \textbf{Comparative Analysis:} Generate difference maps (base - LoRA) to visualize tuning effects
\end{enumerate}

\subsection*{Technical Details}
\begin{itemize}
    \item \textbf{Models:} TinyLlama base vs. LoRA-tuned (adapters/tinyllama-lora-alpaca)
    \item \textbf{Architecture:} 22 layers, 32 heads/layer (704 total heads)
    \item \textbf{Test Cases:} 5 Alpaca prompts spanning short to multi-turn instructions
\end{itemize}

\section*{Results}

\subsection*{Quantitative Metrics}
\begin{table}[H]
\centering
\small
\begin{tabular}{@{}lcc@{}}
\toprule
\textbf{Metric} & \textbf{Base Model} & \textbf{LoRA Model} \\
\midrule
Instruction-following heads & 47 & 52 \\
Clusters identified & 8 & 9 \\
Average attention entropy & 3.24 & 3.18 \\
\bottomrule
\end{tabular}
\caption{Attention behavior comparison}
\end{table}

\subsection*{Key Findings}
\begin{itemize}
    \item \textbf{LoRA increases specialization:} Fine-tuned model has 5 additional instruction-following heads (52 vs 47)
    \item \textbf{Reduced entropy:} LoRA attention is more focused (3.18 vs 3.24), indicating sharper token attribution
    \item \textbf{Late-layer changes:} Largest attention differences in layers 18-21, where task-specific reasoning emerges
\end{itemize}

\subsection*{Visual Evidence}
\begin{figure}[H]
\centering
\includegraphics[width=0.44\textwidth]{../results/task3/head_clustering_Base.png}
\includegraphics[width=0.44\textwidth]{../results/task3/head_clustering_LoRA.png}\\[0.1cm]
\includegraphics[width=0.70\textwidth]{../results/task3/base_vs_lora_attention_diff.png}
\caption{Top: Clustering dendrograms (base left, LoRA right). Bottom: Attention difference heatmap.}
\end{figure}

\section*{Conclusion}
Attention visualization reveals that LoRA fine-tuning \textbf{refines existing attention patterns} rather than creating entirely new mechanisms. The model develops more \textbf{focused, task-specific attention} with increased head specialization. For practitioners, monitoring late-layer attention serves as a quality indicator during fine-tuning.

\end{document}
